% Slide 2: Project Context & Problem Statement
\begin{frame}{Project Context}
    \begin{itemize}
        \item \textbf{Objective:} Simulate a dynamic environment where periodic tasks are activated/deactivated via a remote supervisor (Client) through a central controller (Server).
        \item \textbf{System Model:}
        \begin{itemize}
            \item \textbf{Server:} Manages a pool of clients and a pool of real-time worker threads.
            \item \textbf{Clients:} Issue commands to modify the system state (Activate, Block, Stop, Quit).
            \item \textbf{Tasks:} Periodic routines with predefined execution times and periods.
        \end{itemize}
        \item \textbf{The Challenge:} Managing high-frequency state changes across multiple threads (Connection Handlers, Task Workers, and the Main Server Loop) without incurring race conditions or deadlocks.
    \end{itemize}
\end{frame}

% Slide 3: Architectural Overview
\begin{frame}{System Architecture}
    \begin{center}
        \begin{tikzpicture}[node distance=2cm]
            \node (client) [draw, rectangle, fill=blue!10] {Client (Command Generator)};
            \node (server) [draw, rectangle, fill=orange!10, below of=client, node distance=3cm] {Server (Execution Supervisor)};
            \node (tasks) [draw, cylinder, fill=green!10, below of=server, node distance=3cm] {Active Task Pool (Worker Threads)};
            
            \draw[->, thick] (client) -- node[right] {TCP/IP Command} (server);
            \draw[->, thick] (server) -- node[right] {Status/Answers} (client);
            \draw[->, thick] (server) -- node[right] {Spawn/Control} (tasks);
        \end{tikzpicture}
    \end{center}
    \begin{itemize}
        \item \textbf{Decoupled Logic:} The server separates connection management from task execution.
        \item \textbf{Concurrency Model:} One thread per client (Connection Handler) and one thread per active task (Worker Thread).
    \end{itemize}
\end{frame}

% Slide 4: Data Structure Analysis
\begin{frame}{Core Data Structures}
    To ensure thread safety, the system relies on two primary shared arrays:
    \begin{enumerate}
        \item \texttt{Client clients[MAX\_THREADS]}:
        \begin{itemize}
            \item Tracks connected clients, their socket FDs, and assigned ports.
            \item \textbf{Protected by:} \texttt{mutex\_clients\_array}.
        \end{itemize}
        \item \texttt{ActiveTask tasks\_active[MAX\_NUMBER\_ACTIVE\_TASKS]}:
        \begin{itemize}
            \item Tracks currently running periodic tasks, their IDs, and thread handles.
            \item \textbf{Protected by:} \texttt{active\_tasks\_mutex}.
        \end{itemize}
    \end{enumerate}
    \vspace{0.3cm}
    \textit{Note: These structures represent the "Critical Sections" where race conditions are most likely to occur during simultaneous command processing.}
\end{frame}

% Slide 5: Command Logic: Activation
\begin{frame}{Command Logic: \texttt{ACTIVATE} Task}
    When a client sends an \texttt{ACTIVATE <ID>} command, the Server follows this protocol:
    \begin{enumerate}
        \item \textbf{Parsing:} Extract task ID from the command string.
        \item \textbf{Schedulability Check:} Run \texttt{is\_schedulable()}. This is a read-only operation on the task array.
        \item \textbf{Resource Acquisition:} 
        \begin{itemize}
            \item Attempt to decrement the \texttt{free\_slots\_sem}.
            \item Lock \texttt{active\_tasks\_mutex}.
            \item Locate an empty slot in \texttt{tasks\_active}.
        \end{itemize}
        \item \textbf{Thread Spawning:} Call \texttt{pthread\_create} for \texttt{handling\_active\_task}.
        \item \textbf{Unlock:} Release \texttt{active\_tasks\_mutex}.
    \end{enumerate}
\end{frame}

% Slide 6: Command Logic: Deactivation & Termination
\begin{frame}{Command Logic: \texttt{BLOCK}, \texttt{QUIT}, and \texttt{STOP}}
    \begin{itemize}
        \item \textbf{\texttt{BLOCK} (Deactivate Task):}
        \begin{itemize}
            \item Finds the specific instance in \texttt{tasks\_active}.
            \item Sets \texttt{active = 0} inside a \texttt{active\_tasks\_mutex} block.
            \item \textbf{Crucial:} The worker thread, seeing \texttt{active == 0}, terminates its loop.
            \item Calls \texttt{sem\_post} to signal a slot is now free.
        \end{itemize}
        \item \textbf{\texttt{QUIT} (Client Disconnect):}
        \begin{itemize}
            \item Closes the socket and calls \texttt{rmvclient()}.
            \item Uses \texttt{mutex\_clients\_array} to update the client buffer.
        \end{itemize}
        \item \textbf{\texttt{STOP} (Global Shutdown):}
        \begin{itemize}
            \item Updates the global \texttt{close\_server} flag.
            \item \textbf{Synchronization:} Uses \texttt{mutex\_stopserver} to ensure atomic visibility of the shutdown signal to all threads.
        \end{itemize}
    \end{itemize}
\end{frame}

% Slide 7: The Danger of Race Conditions
\begin{frame}{The Race Condition Problem}
    \begin{alertblock}{Why Synchronization is Mandatory}
        In a multi-threaded environment, the following "Check-then-Act" scenarios are hazardous:
    \end{alertblock}
    \begin{itemize}
        \item \textbf{Scenario A (Task Slot Overlap):} Two clients send \texttt{ACTIVATE} simultaneously. Without \texttt{active\_tasks\_mutex}, both might find the same \texttt{free\_pos}, leading to one task overwriting the other.
        \item \textbf{Scenario B (Server Saturation):} A client attempts to connect when the server is at \texttt{MAX\_THREADS}. Without proper signaling, the server might enter an inconsistent state or crash.
        \item \textbf{Scenario C (Dirty Reads):} A worker thread checking its \texttt{active} status while the Connection Handler is deactivating it.
    \end{itemize}
\end{frame}

% Slide 8: Synchronization Strategy: Mutexes
\begin{frame}{Synchronization Strategy I: Mutexes}
    The implementation utilizes several \texttt{pthread\_mutex\_t} objects to enforce **Mutual Exclusion**:
    \begin{description}
        \item[\texttt{mutex\_clients\_array}] Protects the integrity of the client registry. Ensures only one thread modifies the client list at a time.
        \item[\texttt{active\_tasks\_mutex}] Protects the task registry. Critical for preventing corruption when tasks are added or removed.
        \item[\texttt{mutex\_stopserver}] Protects the global shutdown flag, preventing a race between the "Stop" command and the main loop's exit condition.
        \item[\texttt{log\_mutex}] \textbf{Output Integrity:} Essential for real-time debugging. Prevents multiple threads from interleaving their \texttt{printf} calls, which would render the server logs unreadable.
    \end{description}
\end{frame}

% Slide 9: Synchronization Strategy: Semaphores
\begin{frame}{Synchronization Strategy II: Semaphores}
    The system uses a \textbf{Counting Semaphore} (\texttt{free\_slots\_sem}) to manage the \texttt{tasks\_active} array:
    \begin{itemize}
        \item \textbf{Initialization:} \texttt{sem\_init(\&free\_slots\_sem, 0, MAX\_NUMBER\_ACTIVE\_TASKS)}.
        \item \textbf{Resource Tracking:} The semaphore acts as a counter for available slots.
        \item \textbf{Non-Blocking Allocation:} 
        \begin{itemize}
            \item \texttt{sem\_trywait} is used during activation. 
            \item If the semaphore is 0 (no slots), the server immediately rejects the task with a \texttt{FULL CAPACITY} message instead of hanging.
        \end{itemize}
        \item \textbf{Reclamation:} When a task is de-activated, \texttt{sem\_post} increments the counter, allowing the next \texttt{ACTIVATE} command to proceed.
    \end{itemize}
\end{frame}

% Slide 10: Synchronization Strategy: Condition Variables
\begin{frame}{Synchronization Strategy III: Condition Variables}
    To handle the "Server Full" state for clients, we employ \texttt{pthread\_cond\_t (\texttt{roomAvailable})}:
    \begin{enumerate}
        \item \textbf{The Producer (Server Loop):} When a new connection arrives via \texttt{accept()}, the server attempts to add the client.
        \item \textbf{The Consumer (Connection Handler):} If \texttt{findFreePosition()} returns -1 (Server Full), the thread enters:
        \begin{itemize}
            \item \texttt{pthread_mutex\_lock(\&mutex\_clients\_array)}
            \item \texttt{pthread_cond\_wait(\&roomAvailable, \&mutex\_clients\_array)}
        \end{itemize}
        \item \textbf{The Signal:} When \texttt{rmvclient()} is called (a client disconnects), \texttt{pthread\_cond\_signal(\&roomAvailable)} is invoked, waking up a waiting thread to occupy the new slot.
    \end{enumerate}
\end{frame}

% Slide 11: Summary of Synchronization Primitives
\begin{frame}{Summary of Synchronization Primitives}
    \begin{table}
        \centering
        \small
        \begin{tabular}{|l|l|l|}
        \hline
        \textbf{Primitive} & \textbf{Variable} & \textbf{Purpose} \\ \hline
        Mutex & \texttt{mutex\_clients\_array} & Client array integrity \\ \hline
        Mutex & \texttt{active\_tasks\_mutex} & Task array integrity \\ \hline
        Mutex & \texttt{log\_mutex} & Thread-safe I/O (Console) \\ \hline
        Mutex & \texttt{mutex\_stopserver} & Atomic shutdown flag \\ \hline
        Semaphore & \texttt{free\_slots\_sem} & Counting available task slots \\ \hline
        Cond Var & \texttt{roomAvailable} & Managing client connection overflow \\ \hline
        \end{tabular}
        \caption{Mapping of primitives to system constraints.}
    \end{table}
\end{frame}

% Slide 12: Conclusion
\begin{frame}{Conclusion}
    \begin{itemize}
        \item \textbf{Robustness:} The system successfully bridges the gap between high-level TCP commands and low-level real-time thread management.
        \item \textbf{Concurrency Safety:} By utilizing a tiered synchronization approach (Mutexes for state, Semaphores for resources, and Condition Variables for flow control), the system avoids the classic pitfalls of multi-threaded C programming.
        \item \textbf{Scalability:} The architecture allows for $N$ clients and $M$ tasks to operate simultaneously, governed strictly by the defined synchronization boundaries.
    \end{itemize}
    \vspace{0.5cm}
    \centering
    \textbf{\Large Questions?}
\end{frame}
